%=====================================================================
% 第 1 节  引言
% 原文：Section 1. Introduction（arXiv:1702.06499v2, p. 1–2）
%=====================================================================
\section{引言}

\begin{keytakeaways}
\begin{itemize}
  \item \textbf{物理图像：}超对称规范场论与共形场论（CFT）之间存在一系列引人注目的显式联系。Alday--Gaiotto--Tachikawa 发现了其中最为显式的一类（AGT 对应），它与 Nekrasov 及其合作者命名为 BPS/CFT 对应的另一族结果相互交叠；更早地，Kapustin 与 Witten 揭示了 $N=4$ 超 Yang--Mills 理论与几何 Langlands 对应之间的联系，这一发展也自然地融入正在成形的图景。
  \item \textbf{数学工具：}通向几何 Langlands 对应的两条构造进路——Beilinson--Drinfeld 基于共形场论的进路与 Kapustin--Witten 基于四维超对称规范场论的进路；正文还将自由使用关于 Hitchin 模空间的标准定义与结果。
  \item \textbf{关键结论：}上述三条线索（Beilinson--Drinfeld 进路、Kapustin--Witten 进路、AGT 对应）之间的关系在许多方面仍不完整；本文所报告的结果为澄清这些关系带来若干启示，且限于底李代数为 $\mathfrak{sl}_2$ 的情形，更多细节与高秩情形留待后续系列论文。
  \item \textbf{阅读约定：}正文自由使用关于 Hitchin 模空间的标准定义与结果，其简要概述以术语表形式收于附录 \ref{Gloss}；术语首次出现时以上标标记，如 term$^{g)}$。
\end{itemize}
\end{keytakeaways}

\bigskip
\noindent\rule{\textwidth}{0.6pt}
\bigskip

过去几年中，人们发现了超对称规范场论与共形场论（conformal field theory, CFT）之间若干引人注目的联系。其中最为显式的一类由 Alday、Gaiotto 与 Tachikawa 发现\cite{AGT}，如今常被称为 AGT 对应。它与另一族结果相互交叠，Nekrasov 及其合作者将后者命名为 BPS/CFT 对应；关于这一主题的系列论文中的第一篇，以及这一方向先前工作的文献指引，见 \cite{N15}。一个更早的发展——Kapustin 与 Witten 所揭示的 $N=4$ 超 Yang--Mills 理论（SYM）与几何 Langlands 对应之间的关系\cite{KW}——则自然地融入这一正在成形的图景\cite{NW,NRS}。不过可以指出，Beilinson 与 Drinfeld 处理几何 Langlands 对应的进路与 CFT 有密切关联（综述及进一步文献见 \cite{Fr07}），而这些关联在 Kapustin 与 Witten 的进路中并不显然。上述主题之间的若干联系已在 \cite{NW,T10,NRS,GW} 中有所讨论；近期工作 \cite{KP,AFO} 表明，至少其中一部分关系容许进一步的形变，其动因来自超对称规范场论与弦论。然而，这一图景在许多方面似乎仍不完整。

本文是为 String-Math 2016 会议文集准备的，报告了一些结果，它们为澄清 Beilinson 与 Drinfeld 的进路、Kapustin 与 Witten 的进路以及 AGT 对应这三者之间的关系带来若干启示。为使本文篇幅保持在合理范围内，我们将主要限于讨论底李代数为 $\mathfrak{sl}_2$ 的情形。这里所述的结果是作者与 Ioana Coman-Lohi 共同推进的一个更大项目的一部分；一系列论文正在准备之中，将包含更多细节、对数学方面的扩展讨论以及对高秩情形的讨论。

在正文中，我们将自由使用若干关于 Hitchin 模空间的标准定义与结果。附录 \ref{Gloss} 以术语表的形式汇集了一份非常简短的概述。若某一术语在术语表中有对应条目，则该术语首次出现时将带上标，如 term$^{g)}$。

\bigskip
\noindent\rule{\textwidth}{0.6pt}

% [Tutor Note: 记号提醒：原文此处写 $N=4$（非 $\mathcal N=4$），而第 2.2 节又出现 $\mathcal N=(4,4)$，两种写法并存是原文自身的不统一；译文逐字照录、未作统一，引用时不要“顺手”改成 $\mathcal N=4$。另按术语表，“super-Yang--Mills theory (SYM)”统一作“超 Yang--Mills 理论”，连字符是 en-dash（--），不是单个“-”。]
% [Tutor Note: 含糊点拨：原文只说“certain $N=2$ supersymmetric gauge theories”，未指明具体理论。据 \cite{AGT} 的原始陈述（译者补充，非原文内容）：$SU(2)$、$N_f=4$ 的 $\mathcal N=2$ 超共形 QCD 在 $\Om$-背景（参数 $\ep_1,\ep_2$）下的瞬子配分函数，等于球面上插入四个 Liouville 顶点算符的 Virasoro 共形块；中心荷 $c=1+6(b+b^{-1})^2$，其中 $b^2=\ep_1/\ep_2$（依约定）。Nekrasov--Witten 对此的场论解释见本文第 2.3 节。]
% [Tutor Note: 跳步点拨：“Beilinson--Drinfeld 的进路与 CFT 有密切关联，而这些关联在 Kapustin--Witten 的进路中并不显然”——被省略的机制是：BD 构造把几何 Langlands 对应中的 $\CD$-模实现为临界水平 $k=-h^\vee$ 处仿射 Kac--Moody 代数 $\hat{\mathfrak g}_k$ 的共形块，Ward 恒等式赋予共形块层以 $\CD$-模结构；该构造之所以可能，关键在于 Feigin--Frenkel 定理给出 $\CU(\hat{\mathfrak g}_{-h^\vee})$ 的中心同构于形式圆盘上 ${}^L\mathfrak g$-opers 的空间，从而本征值 $E_i$ 可用来参数化 opers。详见本文第 2.1 节及 \cite{Fr07,FF92}。]
% [Tutor Note: 文献提示：\cite{KP,AFO} 所说的“further deformation”具体指：Kimura--Pestun 的 quiver $W$-代数（把 BPS/CFT 一侧的 $W$-代数结构推广到 $\Om$-形变情形）与 Aganagic--Frenkel--Okounkov 的量子 $q$-Langlands 对应（把几何 Langlands 对应沿 $q$-形变方向推广）。原文正文未展开，此处系译者据参考文献条目补出，供核对。]
% [Tutor Note: 本节小结——引言只交代动机与范围，技术内容自第 2 节开始：2.1 为 BD/CFT 进路，2.2 为 KW/规范场论进路，2.3 为 NW/AGT，2.4 指出前两条进路的共同根源（六维 $(2,0)$ 理论在 $\CM^6=\CM^4\times C$ 上的两种紧化，但所用的拓扑扭转不同），2.5 提及 Betti 几何 Langlands 与 Cordova--Jafferis 两条旁支。另：附录 \ref{Gloss} 的 \label{Gloss} 定义在 sec8.tex 中，单独编译本节时该 \ref 会显示为 ??，须整篇编译方可解析，属预期行为。]
